\documentclass[11pt]{article}

% --------------------
% Packages
% --------------------
\usepackage[a4paper,margin=1in]{geometry}
\usepackage{setspace}
\usepackage{enumitem}
\usepackage{hyperref}
\usepackage{titlesec}
\usepackage{amsmath, amssymb}

\hypersetup{
    colorlinks=true,
    linkcolor=blue,
    urlcolor=blue
}

\onehalfspacing

% Section formatting
\titleformat{\section}{\large\bfseries}{\thesection.}{0.5em}{}
\titleformat{\subsection}{\normalsize\bfseries}{\thesubsection}{0.5em}{}

% --------------------
% Title
% --------------------
\title{\textbf{Lecture 3 --- CSE400: Fundamentals of Probability in Computing}\\
\vspace{0.3em}
\large Introduction to Probability Theory (Course Orientation)}
\author{}
\date{}

\begin{document}
\maketitle

% --------------------
\section{Course Introduction}

\begin{itemize}[leftmargin=1.5em]
    \item \textbf{Course:} CSE400 -- Fundamentals of Probability in Computing
    \item \textbf{Lecture:} 3 (Orientation)
    \item \textbf{Instructor:} Dr. Dhaval Patel, PhD (Associate Professor)
    \item \textbf{Affiliation:} SEAS -- Ahmedabad University, MICxN Research Lab
    \item \textbf{Date:} January 13, 2026
\end{itemize}

% --------------------
\section{Learning Mindset}

\subsection*{Growth Mindset}
\begin{itemize}
    \item Failure is an opportunity to grow.
    \item I like to try new things.
    \item I can learn to do anything I want.
    \item Challenges help me grow.
    \item My effort and attitude determine my abilities.
    \item Feedback is constructive.
    \item I am inspired by the success of others.
\end{itemize}

\subsection*{Fixed Mindset}
\begin{itemize}
    \item Failure is the limit of my abilities.
    \item I’m either good at it or I’m not.
    \item My abilities are unchanging.
    \item I don’t like to be challenged.
    \item My potential is predetermined.
    \item When I’m frustrated, I give up.
    \item I stick to what I know.
\end{itemize}

% --------------------
\section{Course Outline}

\subsection{CSE400: General Course Information}
\begin{itemize}
    \item Team
    \item Active Learning Platform: \textbf{Campuswire}
    \item Schedule
    \item Grading
\end{itemize}

\subsection{Why Should We Learn CSE400?}
\begin{itemize}
    \item Example: Daily life conversations
\end{itemize}

\subsection{Engineering Applications}
\begin{itemize}
    \item Speech Recognition
    \item System Radar System
    \item Communication Network
\end{itemize}

% --------------------
\section{Course Team}

\subsection*{Instructor}
\begin{itemize}
    \item Dr. Dhaval Patel
    \item Role: Instructor
    \item Place: Faculty Office (Room--210)
    \item Areas of Interest: xG Networks, Applied ML/DL/RL/AutoML, Intelligent Transportation Systems, Life Sciences, Behaviour Modelling using AI
    \item Email: \href{mailto:dhaval.patel@ahduni.edu.in}{dhaval.patel@ahduni.edu.in}
\end{itemize}

\subsection*{Student Team (TAs / Project Members)}
\begin{itemize}
    \item Deep Patel --- Reinforcement Learning and Pinching Antenna Systems
    \item Prapti Patel --- Smart Sensing Framework using Kolmogorov--Arnold Networks for Detection in 5G Environments; FAN-F for Spectrum Sensing
    \item Raj Koticha --- Proactive Self-Aware Multi-Agent Reinforcement Learning for Resource Management
    \item Ritu Patel --- Intelligent Transportation System
    \item Rushi Moliya --- Optimizing UAV Deployment and Multi-UAV 3D Deployment
    \item Ura Modi --- Pinching Antenna Systems
\end{itemize}

% --------------------
\section{Active Learning and Class Discussion}

\textbf{Course Website:} Campuswire (two sections provided).

\subsection*{Why Use Campuswire?}
\begin{itemize}
    \item Builds confidence via anonymous participation and back-channel communication.
    \item Enables collaborative and active learning.
    \item Supports real-time feedback using polling.
    \item Allows personal communication with instructor/TAs via direct messages.
\end{itemize}

% --------------------
\section{Schedule}

\subsection*{Lecture Sessions}
\begin{itemize}
    \item Section--1: 9:30 AM -- 11:00 AM (Tuesday, Thursday), GICT Room--136
    \item Section--2: 1:00 PM -- 2:30 PM (Tuesday, Thursday), GICT Room--137
\end{itemize}

\subsection*{TA Hours}
In person / online; timings to be finalized.

% --------------------
\section{Connecting with the Instructor}

\begin{itemize}
    \item Contact hours: 24$\times$7 through Campuswire.
    \item Best practice: post queries on Campuswire.
    \item Direct messaging for private discussions.
    \item External engagement: UGRP-8 (2026), offline projects, counseling, or short meetings (via office email).
\end{itemize}

\textbf{Important Note:} LaTeX tutorial and assignment submission guideline.

% --------------------
\section{Project Component}

\begin{itemize}
    \item \textbf{Project Weightage:} 30\%
\end{itemize}

\subsection*{Project Team Formation}
Deadline: 17 January 2026 (Saturday, End of Day).

% --------------------
\section{Project Execution Guidelines}

\subsection*{Major Milestones (M1--M6)}

\begin{itemize}
    \item \textbf{M1 -- Kickstart:} Team formation, area identification, background, motivation, problem formulation.
    \item \textbf{M2 -- Mathematical Modeling:} Random variables, PMF/PDF, CDF, multivariate RVs, joint PMF/PDF/CDF.
    \item \textbf{M3 -- Coding:} Simulation and computation.
    \item \textbf{M4 -- Inference:} Choose, understand, and code a randomized algorithm.
    \item \textbf{M5 -- Randomized Algorithms:} Apply randomized algorithms to the domain problem and compare with deterministic approaches.
    \item \textbf{M6 -- Derive Bounds and Analysis:} Compile and submit final deliverables.
\end{itemize}

\subsection*{Evaluation}
\begin{itemize}
    \item Deliverables: Codes, reports, videos, etc.
    \item Assessment: Before and after mid-semester.
    \item Final evaluation: Project viva and submission.
\end{itemize}

% --------------------
\section{Submissions}

\subsection*{Submission \#1 --- Concept Evolution Maps}
Tools: Miro; diagrams.net (draw.io).

\subsection*{Submission \#2 --- Scribe: Learning Reflection Logs}
\begin{itemize}
    \item Types: (a) Lecture Scribes, (b) Project Scribe
    \item Two groups per lecture generate scribes.
    \item Minimum length: 8--10 pages.
    \item Includes process and decision-making logs: constraints, alternatives, final decisions, evidence, and trade-off matrices.
    \item Total of 6 bi-weekly submissions.
\end{itemize}

\subsection*{Submission \#3 --- Multimodal Artifacts}
\begin{itemize}
    \item One video per milestone (approximately 10--15 minutes).
    \item Explanation of milestones and coding/simulation.
    \item Think-aloud videos, one-minute insight videos, or project demos.
\end{itemize}

% --------------------
\section{Undergraduate Research Programme (UGRP)}

\begin{itemize}
    \item Multidisciplinary learning (Arts, Science, Management).
    \item Experiential and research-driven approach.
    \item 4D Model: \textbf{Discover + Design + Develop + Deliver}.
    \item T-shaped individual: depth in one discipline with breadth across disciplines.
\end{itemize}

Slides also present MICxN lab activities, industry projects, international collaborations, alumni outcomes, conference publications, awards, and UGRP scholar case studies.

% --------------------
\section{Conclusion}

The lecture concludes with an open Q\&A session and instructor contact information.

\vspace{1em}
\noindent\textit{Source: Lecture 3 slide deck (uploaded PDF)}

\end{document}
